% Solutions of Book 4, chapter 5 (cpp-standard-library).
\section*{Chapter \ref{ch:b2:cpp-standard-library} --- The Standard Library: Containers, Iterators and Algorithms}

\begin{solution}{exo:b2:cpp-standard-library:1}
\omoutput{code/bachelor-2/05-cpp-standard-library/outputs/exo_algorithms.txt}{The six lines.}
a: three 4s. b: \code{max_element} returns the \emph{first} largest, at position 0. c: the type of the result is the
type of the initial value: with \code{0}, an \code{int}, each partial sum $0 + 0.5$ is converted back to
\code{int}, 0; with \code{0.0} the sum is 2. d: \code{unique} removes only \emph{consecutive} duplicates, leaving
1 2 1 3 in the first four places. e: \code{remove} shifts the kept elements to the front and returns the end of
them, 3, but the size is still 6: an algorithm cannot change a \omterm{def:b2:cpp-standard-library:container}{container}'s size, so the elements after the
returned position remain, with unspecified values. f: after \code{erase} the size is 3 (1 2 3, in order), and
\code{binary_search} finds 3. Line e answers the question.
\end{solution}

\begin{solution}{exo:b2:cpp-standard-library:2}
After erasing at position $i$, the next element slides into position $i$, and \code{i++} skips it: 12 follows 8
and 6 follows 6, so both survive, and 10, the last, after 7, is never reached as the size shrinks. Fix 1: increment
\code{i} only when nothing was erased (\code{if (even) erase; else i++;}), still $O(n^2)$ moves. Fix 2, linear:
\code{std::erase_if(v, is_even);}, with \code{is_even} a lambda testing the remainder modulo 2, which is \code{remove_if} followed by one \code{erase}.
\end{solution}

\begin{solution}{exo:b2:cpp-standard-library:3}
\omcode{code/bachelor-2/05-cpp-standard-library/exercises/top_words.hpp}{9}{21}{The $k$ most frequent words}
\code{partial_sort} makes about $m \log_2 k$ comparisons (it keeps a heap of the $k$ best), against $m \log_2 m$ for a
full sort; for the 33 words of the town and $k = 3$, about 52 against 166. The comparator breaks ties by the word,
so the result does not depend on the algorithm's internal order. \code{tests/word_count_test.cpp} checks that the
first $k$ agree with \code{by_frequency} on the town's titles.
\end{solution}

\begin{solution}{exo:b2:cpp-standard-library:4}
The map: the $i$-th insertion searches a tree of $i - 1$ keys, about $\log_2 i$ comparisons, so about
$\sum_{i=1}^{n} \log_2 i = \log_2 n! \approx n \log_2 n - 1.44\,n$ comparisons, plus $n$ allocations and
rebalancings of constant amortised cost. The hash table: one hash per insertion, $n$ in all, plus the rehashes. If
the bucket count doubles each time the load factor passes 1, the rehashes move about
$n/2 + n/4 + \dots < n$ elements; the libraries keep the hash in the node or recompute it, at most $n$ more hashes.
Total: $O(n)$ on average, against $O(n \log n)$; the measured lookups of
\cref{out:b2:cpp-standard-library:lookups} show the gap at large $n$.
\end{solution}

\begin{solution}{exo:b2:cpp-standard-library:5}
\omoutput{code/bachelor-2/05-cpp-standard-library/outputs/exo_captures.txt}{The four lines.}
Line 1: \code{by_value} copied \code{fine} when the lambda was created, 10; \code{by_ref} reads the variable, now
25. Line 2: \code{copy} is a copy of the object \code{counter} taken while its \code{n} was 0; three calls of
\code{counter} give 3, the first call of \code{copy} 1. Line 3: \code{[=]} copied \code{fine} (25) at creation, 50;
\code{[&]} reads it after it was set to 0. Line 4: the lambda object stores its \omterm{def:b2:cpp-standard-library:lambda}{captures}: one \code{int}, 4 bytes,
for \code{by_value}, and one reference, stored as an address, 8 bytes, for \code{all_ref}.
\end{solution}

\begin{solution}{exo:b2:cpp-standard-library:6}
The lambda \omterm{def:b2:cpp-standard-library:lambda}{captures} \code{p}, a local variable of \code{starts_with_street}, by reference; the function returns,
its frame is gone, and \code{p} with it. The call in \code{main} reads \code{p.size()} through the dangling reference.
Fix: capture by value, \code{[p]}, or move it in, \code{[p = std::move(p)]}: the string then lives in the lambda
object, which \code{std::function} owns. With the default options, the sanitizer does not watch a returned frame,
whose stack memory is soon reused by the next call; the option makes it place frames in separately allocated
memory that it marks as dead after the return. A freed heap block, as in \cref{out:b2:cpp-standard-library:invalidation},
is always watched.
\end{solution}

\begin{solution}{exo:b2:cpp-standard-library:7}
\omcode{code/bachelor-2/05-cpp-standard-library/exercises/by_month.hpp}{9}{25}{Member, then month}
\code{month_of} inverts \code{day_number}. A \code{map} visits members and months in increasing order, which is
how the table is read and printed; an \code{unordered_map} would need sorting first, and for 60 members and 12
months the speed of lookups does not matter. \code{tests/loans_test.cpp} checks that each member's months add up
to the member's loans, and that the whole adds up to 400.
\end{solution}

\begin{solution}{exo:b2:cpp-standard-library:8}
(a) Let $k = \lfloor \log_2 n \rfloor$, so $2^k \leq n < 2^{k+1}$, hence $2^k < n + 1 \leq 2^{k+1}$ and
$\lceil \log_2 (n + 1) \rceil = k + 1$. (b) The comparisons a search makes, each with two outcomes, form a binary
decision tree whose leaves are the returned positions. Every one of the $n + 1$ positions (0 to $n$) is the right
answer for some key, so the tree has at least $n + 1$ leaves; a binary tree of height $h$ has at most $2^h$ leaves,
so $2^h \geq n + 1$ and $h \geq \lceil \log_2(n + 1) \rceil$: some key takes that many comparisons. By (a), it is
the bound of \cref{prop:b2:cpp-standard-library:lower-bound}, which no comparison-based search can beat.
\end{solution}

\begin{solution}{exo:b2:cpp-standard-library:9}
\omcodes{code/bachelor-2/05-cpp-standard-library/exercises/lru.hpp}{9-24,25-44}{An LRU cache}
\code{splice} moves a node from one place of a list to another without copying it or invalidating any iterator,
and erasing the back node \omterm{def:b2:cpp-standard-library:invalidation}{invalidates} only that node's iterators, which the map entry erased with it was the only
one to hold. The iterators stored in the map therefore stay valid for as long as their entries exist. With a
\code{vector} instead of the list, every move to the front would shift the others and invalidate their stored
positions. \code{tests/lru_test.cpp} compares 15\,000 random operations with a vector kept in recency order.
\end{solution}

\begin{solution}{exo:b2:cpp-standard-library:10}
\omcode{code/bachelor-2/05-cpp-standard-library/exercises/interval_map.hpp}{12}{48}{An interval map on
\code{std::map}}
\code{upper_bound(key)} is the first change after \code{key}, so the one before it, if any, is the last change at or
before \code{key}. \code{assign} saves the value that must resume at \code{hi}, erases the changes inside
$[lo, hi)$, writes the change at \code{lo} and the change back at \code{hi}, then removes either of them if it
repeats the value before it. It costs $O(\log n)$ plus the number of erased changes, each of which was inserted
once. The test checks every key from $-10$ to 49 and the number of stored changes after each of 3000 random
assignments.
\end{solution}

\begin{solution}{exo:b2:cpp-standard-library:11}
(a) The printed points are 1, 2, 3, 5, 9, \dots, 513: the capacity doubles, starting at 1. For $n \geq 1$
\code{push_back}s there are $1 + \lceil \log_2 n \rceil$ reallocations (11 for $n = 1000$), and the moves add up to
$1 + 2 + 4 + \dots + 2^{\lceil \log_2 n \rceil - 1} < 2n$. (b) Exactly when $i = 1$ or $i - 1$ is a power of two: the
$i$-th call finds $i - 1$ elements in a full block. (c) \code{due.reserve(8);} before the loop: no
\code{push_back} reallocates, and the end iterator held by the range-\code{for} still marks four elements, so the
loop reads 3, 8, 12 and 20 and stops. It is fragile: one more renewal than reserved and the bug returns, silently.
\end{solution}

\begin{solution}{exo:b2:cpp-standard-library:12}
\omcode{code/bachelor-2/05-cpp-standard-library/exercises/workload.hpp}{9}{40}{One workload, three queues}
\omoutput{code/bachelor-2/05-cpp-standard-library/outputs/exo_workload.txt}{One run on the author's laptop
(\code{exo_workload.cpp}).}
By the guarantees, the heap and the multiset take $O(\log w)$ per operation and the sorted vector $O(w)$ per
insertion (the smallest is at the back, so serving is $O(1)$). The measurement agrees for large $w$: at 100\,000
waiting, the vector is about forty times slower than the heap. For a few hundred waiting, the sorted vector is the
fastest, and the multiset is slower than the heap at every size (a node allocation per insertion). Choice: the heap,
which is never far from the best; the sorted vector only if the list is known to stay short. The test in
\code{tests/exercises_test.cpp} checks the order of service, not the times.
\end{solution}

\begin{solution}{pb:b2:cpp-standard-library:1}
\textbf{1.} \code{sort_by_due} in \code{exercises/loan_desk.hpp}: compare \code{due}, then \code{id}.

\textbf{2.} \code{lower_bound(d)} is the first loan due on or after $d$, \code{upper_bound(d)} the first due after
$d$; between them, the loans due on $d$. \code{equal_range} returns both as a pair, in one call. The comparator must
compare a loan with a day in both orders, as \code{DueBefore} does.

\textbf{3.} At most $\lfloor \log_2 400 \rfloor + 1 = 9$ each.

\textbf{4.} \code{due_between(sorted, d, d + 7)}: from \code{lower_bound(d)} to \code{lower_bound(d + 7)}. Consecutive
weeks $[d, d + 7)$ and $[d + 7, d + 14)$ then share no loan and miss none. The test compares every such range with a
count by hand.

\textbf{5.} At \code{upper_bound} of its due date, found in 9 comparisons, but the insertion shifts every later loan:
up to 400 moves.

\textbf{6.} If the vector had spare capacity, the iterators before the insertion point; otherwise none.

\textbf{7.} \code{operator[]} inserts the key with an empty vector when it is missing, which is what grouping needs.

\textbf{8.} \code{operator[]} may insert, so it does not exist on a \code{const} map. \code{at} throws if the key is
missing; \code{find} returns \code{end()}.

\textbf{9.} The counts are the same; the ties of \code{stable_sort} come out in a different order, since the
pairs are now built in the hash table's order instead of by increasing member number.

\textbf{10.} \code{++borrowed[l.book_id]} for each loan, then \code{max_element} with a lambda comparing the counts:
book 45, borrowed 10 times.

\textbf{11.} Loop with \code{auto it = m.begin(); it != m.end();} and no increment in the header; in the
body, \code{it = m.erase(it);} if the member has one loan, \code{++it;} otherwise. \code{erase} returns the iterator
after the erased one, so the erased one is never incremented. Three members are removed. In C++20,
\code{std::erase_if(m, pred)} does the same.

\textbf{12.} Yes for the map: insertions never move nodes. Yes for the \code{unordered_map} too: a rehash
\omterm{def:b2:cpp-standard-library:invalidation}{invalidates} iterators, but pointers and references to the elements survive.

\textbf{13.} 109, as in \cref{ch:b2:cpp-polymorphism}.

\textbf{14.} The late loans come first and \code{partition} returns the boundary; the order within each part is
unspecified. \code{stable_partition} keeps it, at a higher cost.

\textbf{15.} \code{std::erase_if(v, [](const Loan &l) { return l.returned >= 0; });} keeps the 22 loans still out.

\textbf{16.} \code{loans | views::filter(late) | views::take(10)}, then \code{views::transform} with a lambda returning
\code{l.id}, with \code{late} a lambda capturing the day.

\textbf{17.} \code{fines_cents} in \code{loan_desk.hpp}: 52\,850 cents, 528.50. \code{accumulate} returns the type of
its initial value: with \code{0}, each sum would be computed in \code{int}, which would suffice here but overflow
for a large library, and amounts in \code{double} would be truncated.

\textbf{18.} \code{tests/loans_test.cpp} counts with a plain loop and compares with \code{late_count}, and does the
same for the fines.

\textbf{19.} The vector: $O(n)$ moves per insertion, against $O(\log n)$ for the map.

\textbf{20.} \omoutput{code/bachelor-2/05-cpp-standard-library/outputs/crossover.txt}{One run on the author's
laptop.}

\textbf{21.} At $n = 1000$, an insertion moves on average 500 integers, 2\,KB of contiguous memory, which the
processor copies at many bytes per cycle; the binary search reads a 4\,KB array that stays in the fastest cache.
The map allocates a node for each key and follows about ten pointers per search.

\textbf{22.} Above about 10\,000 keys, the average shift exceeds 20\,KB per insertion and grows linearly with $n$:
the cost per key doubles when $n$ doubles, as $O(n)$ per insertion predicts.

\textbf{23.} Its $\log_2 n$ steps grow slowly, but its nodes no longer fit in the caches, and each step becomes a
cache miss.

\textbf{24.} \emph{The crossover size at which a sorted vector stops beating a map, on the laptop}: in the printed
run, between 8000 and 16\,000 keys, about $10^4$; below it the vector is faster per key (by about 20\,\% at
1000), above it the map is, by a factor that doubles with $n$. It lies where the linear cost of shifting,
cheap per element because the memory is contiguous, overtakes the logarithmic cost of the map, high per step
because each step follows a pointer.
\end{solution}
